\chapter{Zones, routes and fixed OD demand}
\section{Convert turning counts to OD demand}
For an isolated intersection, approach and exit zones are a transparent starting point. If the west approach contains 500 PCU/h through, 80 PCU/h left and 70 PCU/h right, define three OD movements rather than treating the three turns as alternative routes to the same destination.
\begin{lstlisting}[language=yaml]
routes:
  - id: R_W_E
    origin_zone_id: Z_W
    destination_zone_id: Z_E
    movement_ids: [T_WE]
  - id: R_W_N
    origin_zone_id: Z_W
    destination_zone_id: Z_N
    movement_ids: [L_WN]
  - id: R_W_S
    origin_zone_id: Z_W
    destination_zone_id: Z_S
    movement_ids: [R_WS]

origin_destination_matrix:
  values:
    Z_W:
      Z_E: 500
      Z_N: 80
      Z_S: 70
\end{lstlisting}
Thus the west approach total is $500+80+70=650$ PCU/h.

\section{OD conservation}
For every OD pair $ij$,
\begin{equation}
\sum_{r\in R_{ij}} q_r=Q_{ij}.
\end{equation}
The matrix total is fixed. Assignment changes only the route distribution.

\section{Assignment convergence}
PySig uses a logit target allocation and method of successive averages. A useful representation is
\begin{equation}
w_r=\exp[-\theta(c_r-c_{\min})]
\end{equation}
and
\begin{equation}
q_r^n=q_r^{n-1}+\frac{1}{n}(\hat q_r^n-q_r^{n-1}).
\end{equation}
Convergence is assessed using maximum route-flow change. OD conservation is an accounting identity, not a behavioural convergence criterion.
